\documentclass[border=10pt]{standalone}
\usepackage{tikz}
\usetikzlibrary{arrows.meta}
\definecolor{myred}{RGB}{190,60,45}

%======== 整体右移量：和前两张图保持一致，只改这一个数即可 ========
\newcommand{\xoff}{2.5}

%======== 损失函数（定义不变，画图时横坐标统一 +\xoff）========
\pgfmathdeclarefunction{L}{1}{%
  \pgfmathparse{0.01*(#1-4)^4 + 1.4*exp(-(#1-4)*(#1-4)/3) - 0.15*(#1-4)}}

\begin{document}
\begin{tikzpicture}[x=0.6cm, y=0.8cm, >=stealth, line width=0.8pt, font=\large,
    garr/.style={myred, ->, >={Stealth[length=5pt, width=4pt]}, line width=0.8pt}]

  % 坐标轴（x 轴左端只留一点尾巴，曲线不再穿过 y 轴）
  \draw[->] (-1.2,0) -- (13.2,0) node[below right] {$\theta$};
  \draw[->] (0,-0.6) -- (0,7.9)  node[above right] {$L(\theta)$};

  % 损失曲线（黑色）
  \draw[black, line width=1pt, domain=-1:9.3, samples=200]
        plot ({\xoff+\x}, {L(\x)});

  % 两个极小：黑点 + 虚线垂直落到 x 轴 + 轴上小刻度 + 轴下标签
  % （线的样式与 lr_small 图完全一致；字号 footnotesize 也与它一致）
  \fill[black] ({\xoff+2},{L(2)})     circle (2.5pt);
  \fill[black] ({\xoff+6.3},{L(6.3)}) circle (2.5pt);

  % 左：局部极小
  \draw[dashed,gray,thin] ({\xoff+2},{L(2)}) -- ({\xoff+2},0);
  \draw[gray,thin]        ({\xoff+2},0.12)  -- ({\xoff+2},0);
  \node[font=\footnotesize,anchor=north] at ({\xoff+2},-0.15) {local minimum};

  % 右：全局极小
  \draw[dashed,gray,thin] ({\xoff+6.3},{L(6.3)}) -- ({\xoff+6.3},0);
  \draw[gray,thin]        ({\xoff+6.3},0.12)    -- ({\xoff+6.3},0);
  \node[font=\footnotesize,anchor=north] at ({\xoff+6.3},-0.15) {global minimum};

  % ---- 学习率合适：同一起点，8 步翻过小山包，收敛到全局最小 ----
  % 落脚点（9 个点 = 8 步）
  \foreach \t in {0.30,1.25,2.10,2.90,3.65,4.40,5.10,5.75,6.30}
    \fill[myred] ({\xoff+\t},{L(\t)}) circle (2pt);

  % 更新箭头
  \foreach \a/\b in {0.30/1.25, 1.25/2.10, 2.10/2.90, 2.90/3.65,
                     3.65/4.40, 4.40/5.10, 5.10/5.75, 5.75/6.30}
    \draw[garr] ({\xoff+\a},{L(\a)}) -- ({\xoff+\b},{L(\b)});
\end{tikzpicture}
\end{document}